\section{İlgili Çalışmalar}
\label{sec:ilgili}

Bu bölüm, çalışmanın konumlandığı literatürü beş tematik küme altında
incelemektedir. Tarama, hedefli (targeted/scoping) bir kapsam taramasıdır;
PRISMA uyumlu sistematik bir tarama değildir ve her kümede alanın en çok
atıf alan/temel referansları önceliklendirilmiştir.

Küme 1, prosedürel arazide çoklu çözünürlük karıştırma ve LOD (detay
seviyesi) geçişini ele alır. Geometry clipmaps
\citep{losasso2004geometry}, ROAM \citep{duchaineau1997roaming} ve
geomipmapping \citep{deboer2000geomipmapping}, tam olarak bu çalışmanın
iki mesh (küresel + yerel yama) çerçevesinin ait olduğu geleneksel
bilgisayar grafikleri hattını oluşturur; üçü de düşük ve yüksek
çözünürlüklü ağların birleştiği sınırda çatlak/dikiş oluşmasını merkezi
bir mühendislik problemi olarak ele alır. Bu çalışmanın bulgusu, üretim
öncesi Gauss'un sonlu yarıçapta kesildiğinde gerçek bir dikiş bırakması,
bu literatürün yirmi beş yılı aşkın süredir çözmeye çalıştığı aynı sınıf
problemin somut ve ölçülmüş bir örneğidir; smoothstep'in kompakt desteği
ise ROAM ve geoclipmap'in kullandığı sınırda tam sıfıra inen ağırlık
fonksiyonu ilkesiyle örtüşür.

Küme 2, kompakt destekli ve sonsuz destekli karıştırma fonksiyonlarının
süreklilik özelliklerini ele alır. Hardy'nin \citep{hardy1971multiquadric}
topografya için önerdiği multiquadric/Gauss tipi radyal baz fonksiyonu, bu
çalışmadaki üretim öncesi Gauss'un ait olduğu aile gibi sonsuz destekli
olup hiçbir sonlu yarıçapta tam sıfıra inmez; bu, ölçülen kalıntı
davranışıyla (yaklaşık $\exp(-1)\approx\SI{37}{\percent}$) örtüşür.
Wendland \citep{wendland1995piecewise} ise bunun çözümünü formalize eder:
verilen bir düzgünlük derecesi için kompakt destekli, minimal dereceli
pozitif tanımlı radyal fonksiyonlar; smoothstep (kübik Hermite,
\citep{ebert2002texturing}), bu ailenin bir boyutlu, düşük dereceli özel
bir örneğidir. Ohtake ve ark.\ \citep{ohtake2003multiscale}, bu çalışmanın
hibrit yönteminin izlediği stratejiyi tanımlar: kompakt destekli baz
fonksiyonlarla çok ölçekli bir hiyerarşi kurup farklı ölçeklerde farklı
karıştırma davranışı elde etmek.

Küme 3, gezegen/ay inişi simülasyonlarında arazi modelleme ve temas
dinamiğini ele alır. ALHAT programının hedefi
\citep{johnson2008overview,nasajpl_alhat}, yani yüz metre içinde hassas
iniş ve LiDAR tabanlı arazi göreli navigasyon, ve LOLA/SELENE tabanlı
dijital yükseklik modellerinin sunduğu doğruluk düzeyi
\citep{barker2016lunar}, gerçek ay inişi mühendisliğinde arazi
doğruluğunun kozmetik değil doğrudan operasyonel bir gereksinim olduğunu
göstermektedir. Bu, çalışmanın ``neden dikiş/süreklilik önemli''
motivasyonunu güçlendirmektedir: ayak temas bölgesindeki bir süreksizlik,
gerçek ALHAT sınıfı sistemlerde tehlike tespiti ve iniş alanı seçimi
kararlarını etkileyebilecek bir hata sınıfına karşılık gelir. Bu
motivasyon, güncel bir uçuş verisiyle de doğrulanmıştır: JAXA'nın SLIM
görevinin (Ocak 2024) görüş-tabanlı navigasyon ve engel tespiti sonuçları
\citep{ishida2025vision}, inişten \SI{50}{\meter} önce yapılan gerçek
zamanlı engel/arazi tespitinin metre mertebesinde iniş hassasiyeti
sağladığını göstermiş; bu, arazi temsilindeki hataların yalnızca
simülasyonda değil, gerçek uçuşta da doğrudan iniş güvenliğine
yansıdığını somutlaştırmaktadır.

Küme 4, fizik tabanlı RL ortamlarında prosedürel arazi randomizasyonunu
ele alır. Rudin ve ark.\ \citep{rudin2022walking}, IsaacGym üzerinde bu
çalışmadaki terrain havuzuna benzer bir zorluk tabanlı
terrain-curriculum kullanır; bu, çalışmanın curriculum ve terrain-havuzu
tasarımının izole bir örnek olmadığını, alanın yerleşik bir pratiği
olduğunu göstermektedir. Mittal ve ark.\ \citep{mittal2023orbit} (Orbit)
ve NVIDIA \citep{nvidia2025isaaclab} (Isaac Lab), doğrudan bu çalışmanın
kullandığı platformun mimari temelini tanımlarken, Tobin ve ark.\
\citep{tobin2017domain} prosedürel randomizasyonun sim-to-real
transferindeki genel motivasyonu sağlar. Ancak bu kaynakların hiçbiri,
arazi randomizasyonu içindeki çoklu çözünürlük karıştırma fonksiyonunun
kendisini (smoothstep, Gauss veya hibrit) bir tasarım değişkeni olarak
ele almamaktadır; hepsi terrain çeşitliliğine odaklanmakta, blend
şeklinin süreklilik ve öğrenme maliyetine odaklanmamaktadır. Bu çalışma
tam olarak bu boşlukta konumlanmaktadır.

Küme 5, hibrit/bölgesel karıştırma fonksiyonu birleştirmesini (partition
of unity) ele alır. Ohtake ve ark.\ \citep{ohtake2003multiscale}, burada
da merkezi kaynaktır: kompakt destekli baz fonksiyonların çok ölçekli,
partition-of-unity tarzı birleştirilmesi, bu çalışmanın hibrit yönteminin
(temas bölgesinde bir fonksiyon, ötesinde başka bir fonksiyon, aralarında
yumuşak bir switch) genel matematiksel çerçevesidir. Bu çerçeve altında,
hibrit switch'in makine hassasiyetinde yeni bir dikiş yaratmaması
beklenen ve literatürle tutarlı bir sonuçtur: iki C1 düzgün fonksiyonun
C1 düzgün bir ağırlıkla dışbükey birleşimi yine C1 düzgündür.

Özetle, Küme 1 problem sınıfını (çoklu çözünürlük geçişinde
dikiş/süreklilik) sağlarken bu literatür ağırlıklı olarak görsel kaliteye
odaklanmaktadır; bu çalışma aynı problemi fiziksel olarak sorgulanabilir
bir yükseklik alanında (ayak teması) sayısal olarak ölçmektedir. Küme 2,
ampirik bulgunun neden doğru olduğunu elli yılı aşkın bir matematiksel
sonuçla temellendirmektedir; bu bir kodlama hatası değil, radyal baz
fonksiyonlarının bilinen bir yapısal özelliğidir. Küme 3, ``neden
önemli'' sorusuna gerçek dünya mühendislik gerekçesi sağlamaktadır. Küme
4, bu çalışmanın deney platformunun literatürdeki yerini ve doldurduğu
boşluğu netleştirmektedir. Küme 5, önerilen hibrit yöntemi izole bir
mühendislik hilesi olmaktan çıkarıp yerleşik bir matematiksel çerçeveye
oturtmaktadır.
